\documentclass[12pt,a4paper]{article}

% ------------------------------------------------
% PAGE FORMATTING
% ------------------------------------------------
\usepackage[a4paper, margin=1in]{geometry}
\usepackage{setspace}
\onehalfspacing
\usepackage{parskip}

% ------------------------------------------------
% MATHEMATICS
% ------------------------------------------------
\usepackage{amsmath}
\usepackage{amssymb}
\usepackage{amsthm}
\usepackage{mathtools}

% ------------------------------------------------
% GRAPHICS AND FIGURES
% ------------------------------------------------
\usepackage{graphicx}
\usepackage{float}
\usepackage{subcaption}

% ------------------------------------------------
% LISTS
% ------------------------------------------------
\usepackage{enumitem}

% ------------------------------------------------
% TABLES
% ------------------------------------------------
\usepackage{booktabs}
\usepackage{array}

% ------------------------------------------------
% HYPERLINKS
% ------------------------------------------------
\usepackage[colorlinks=true, linkcolor=blue, urlcolor=blue, citecolor=blue]{hyperref}

% ------------------------------------------------
% CUSTOM COMMANDS
% ------------------------------------------------
\newcommand{\PPP}{\text{PPP}}
\newcommand{\E}{\mathbb{E}}
\newcommand{\Prob}{\mathbb{P}}

% ------------------------------------------------
% TITLE INFORMATION
% ------------------------------------------------
\title{CSE400 --- Fundamentals of Probability in Computing\\
Milestone 3 Scribe\\
Handover Probability in Drone Cellular Networks\\
Group - s1\_g14\_net}

\date{\today}

\usepackage{titlesec}
\titleformat{\section}{\large\bfseries}{\thesection}{1em}{}
\setcounter{secnumdepth}{0}

\begin{document}

\maketitle
\thispagestyle{empty}

% ------------------------------------------------
% SCRIBE QUESTION 1
% ------------------------------------------------

\section{Scribe Question 1: System Objective and Deterministic Baseline}

This study aims to examine the rate of handovers that occur within a drone cellular network where drones fly and act as base stations (BS) for stationary ground users. The drone's starting positions are assumed to be organized in accordance with a homogeneous Poisson Point Process (PPP) with density $\lambda_0$ within a BS plane; this distribution has been used extensively within stochastic geometry and is representative of the randomly deployment pattern of many types of BS.

The drones will each fly in straight lines at uniformly random angles, resulting in the drone network changing over time as the drone serving any given user assignment to that user may also eventually change. The possible change of the serving drone is referred to as a handover, and it is defined to occur when the serving drone changes prior to time $t$.

Mathematically, the handover event is defined as

\[
H(t) := \{\exists s < t : \arg\min_i x_i(s) \neq \arg\min_i x_i(t)\}
\]

where $x_i(t)$ represents the distance of drone $i$ from the user at time $t$. The handover probability is then defined as the probability that the first handover occurs before time $t$:

\[
P[H(t)]
\]

This probability is the primary indicator of success measured against in this study.

The deterministic baseline algorithm applied to this approach was a nearest-neighbour based association policy whereby the client is always connected to the nearest drone based base station at time $t$, so the serving drone will be as follows at time $t$:

\[
x^*(t) = \arg \min_i ||x_i(t)||
\]

Let $x_i(t)$ represent the location of drone $i$. When one of the drones moves closer to the currently associated drone than to the one that serves them, a handover takes place.

Cellular network performance models that use this baseline are popular because the adjacent base station provides maximum signal strength based on path-loss characteristics. Handover probability is one of the performance measures that can be calculated analytically using the nearest base station for point-to-point model networks, making it the ideal association method for studying mobility in drone-cellular networks.

% ------------------------------------------------
% SCRIBE QUESTION 2
% ------------------------------------------------

\section{Scribe Question 2: Random Variables in the Simulation}

There are multiple random variables in space, including the number of Drone base System (DBSs), the spatial distribution of DBSs, and the mobility of DBSs that represent the distribution and movement properties of DBSs over time. Each of these random variables captures the unknown nature of the DBS's physical installation (deployment) in the system along with how the DBS's move through the system over time.

The number of drones will also be considered a random variable. The DBS's in the system will begin with a homogeneous Poisson point process with a density of $\lambda_0$ over the entire spatial area. Therefore, the total number of drones in an $R$ radius circular simulation area will be Poisson distributed. So,

\[
N \sim Poisson(\lambda_0 \pi R^2)
\]

The variable $\lambda_0$ depicts the spatial density of drone base stations, which is frequently applied as an analytical modeling method by using stochastic geometry to create a model of a wireless network in which drone base stations are randomly deployed over space.

The starting coordinates of the drones are also randomly assigned. In the numerical experiments described here, the drones were uniformly distributed through the circular area. If we let denote a drone’s polar coordinates $(r,\theta)$

\[
r = R\sqrt{U}, \quad \theta \sim U(0,2\pi)
\]

$U$ is a uniform random variable on the interval $[0,1]$. This gives each of the drones a uniform distribution across the entire area.

Another random variable of importance in the simulation is the direction of movement for each drone. Each drone travels in a straight line at an angle that is randomly selected. The angle is assumed to be uniformly distributed over the range from 0 to 360 degrees.

\[
\theta_d \sim U(0,2\pi)
\]

For the Single Speed Model (SSM), all of the drones are moving through the scenario at the same constant speed $(v)$.

In the Different Speed Model (DSM), the drones can travel at different speeds; for example, the velocity of each drone may vary based on the Rayleigh distribution

\[
f_V(v) = \frac{v}{\sigma^2} e^{-v^2/(2\sigma^2)}
\]

or a uniform distribution over a given interval

\[
V \sim U(v_{\min}, v_{\max})
\]

Using these variables, the position of a drone at time $t$ can be expressed as

\[
x_i(t) = x_i(0) + v_i t \cos(\theta_i)
\]

\[
y_i(t) = y_i(0) + v_i t \sin(\theta_i)
\]

The drone's velocity is denoted by $v_i$ and the drone's direction is denoted by $\theta_i$. However, even if the values of these parameters can change randomly, the connection choice remains fixed; at each timestep, the user connects to whichever drone happened to be closest at that time. If the index of the closest drone changes during the course of the simulation, a handover will have taken place.

% ------------------------------------------------
% SCRIBE QUESTION 3
% ------------------------------------------------

\section{Scribe Question 3: Probabilistic Modeling and Fixed Decision Rule}

The way the DBS is located in relation to other drones creates uncertainty based on the spatial distribution and movement pattern of all DBS.

If we assume a homogeneous Poisson Point Process (PPP) with density $\lambda_0$ for the initial locations of the DBS, the total number of drones located in a geographic ground area will be Poisson-distributed.

\[
N \sim Poisson(\lambda_0 A)
\]

A second source of randomness is where each individual drone will move in a straight line towards its randomly selected displacement direction.

\[
\theta \sim U(0,2\pi)
\]

The movement can vary according to various types of speed models; for example, the single speed model (SSM) or the different speed model (DSM).

Rayleigh distribution:

\[
f_V(v) = \frac{v}{\sigma^2} e^{-v^2/(2\sigma^2)}
\]

Using these assumptions, the position of a drone evolves over time according to

\[
x_i(t) = x_i(0) + v_i t \cos(\theta_i)
\]

\[
y_i(t) = y_i(0) + v_i t \sin(\theta_i)
\]

The user will always connect to the nearest drone base station at time $t$:

\[
x^*(t) = \arg \min_i ||x_i(t)||
\]

A handover occurs whenever the closest drone changes its index.

\[
H(t) := \{\exists s < t : \arg\min_i x_i(s) \neq \arg\min_i x_i(t)\}
\]

The main quantity of interest is the handover probability

\[
P[H(t)]
\]

which indicates the likelihood that there will be an instance when a handoff has taken place prior to $t$.

\newpage

% ------------------------------------------------
% SCRIBE QUESTION 4
% ------------------------------------------------

\section*{Scribe Question 4: Model-Implementation Alignment}

\textbf{Translating Stochastic Geometry into Programmatic Logic}

The theoretical mathematical model changes to the Python and MATLAB scripts is a large reliance on Monte Carlo techniques and deterministic numerical in- tegration. The theoretical spatial arrangement of the drone base stations, which are the starting point of the drone base stations. a homogeneous Poisson Point Process on a 2D plane, defined by model is simulated by. modeling the total number of drones as a Poisson distributed variable and locating the drones at random. formally working with polar coordinates in a discrete simulation radius. The nearest-neighbor association policy is enforced programmatically by defining an exclusion zone. The distance to the first serving drone is identified in the scripts and a bool array is used. mask so that no drones are filtered through that radius, exactly resembling the mathematical. definition of the non-serving drone point process. 

\textbf{Kinematic Updates and Empirical Validation}

The random direction mobility models, and the same speed Model of both models are straight-line mobility models. The different speed model (DSM) and the SSM are applied through discrete kinematic. updates. The simulation chooses random speeds, based on given distributions, including Rayleigh or Uniform, and renews the Cartesian coordinates of each drone at each time. Whenever the recalculated minimum distance is less than it was before, handover events are empirically recorded. below the original serving range. To evaluate the exact handover formulas and lower bounds computationally, the theoretical scripts shift from Monte Carlo to deterministic numerical integration, utilizing libraries like SciPy to evaluate the complex geometric intersection areas and probability density functions.

\textbf{The Law of Large Numbers in Action}

The results of the given simulation are a perfect demonstration of the correspondence between the theoretical. model and the empirical representation by the Law of Large Numbers. In the plot produced by using 1,000 realizations, the generated data points (red circles) are simulated, and hence have. strong noise and variance, commonly jaggedly discontinuous with the smooth analytical curves. Such scattering directly owes to the very high spatial randomness of the network. first release and disorganized movement. However, in the plot scaled to 10,000 as realizations, the empiric data becomes smoothing. The simulating points are tracking the theoretic exact SSM curve and DSM lower limit with high accuracy. This visual convergence certifies that the programmed simulation is an accurate reflection of the underlying. mathematical structure of the paper 

% ------------------------------------------------
% SCRIBE QUESTION 5
% ------------------------------------------------

\section*{Scribe Question 5: Cross-Milestone Consistency and Change}

\textbf{From Abstract Mathematics to Concrete Execution}

The greatest development following the above milestones is the crystallization of abstract probabilistic models into concree program implementation. In the earlier math- The concepts such as the exclusion zone and the inhomogeneous density of are broken up mathematically. non-serving drones were purely theoretical constructions that were determined by the void probability. theorems of formula and displacement. At the present implementation phase, the following notions were present. have now shifted to literal manipulations of arrays and spatial coordinates updates. Similarly The complicated double and triple integrals that have been used to establish the handover probability are intricate. Lower bounds have not just been symbolic representations on paper, but nested numerical. Functions of quadrature determined in the code. 

\textbf{Validating Theoretical Assumptions}

The core theoretical assumptions established in earlier milestones have remained remarkably consistent and are now visually verifiable. The critical observation that the time-correlated, heterogeneous velocities of the Different Speed Model result in a lower handover probability than the Same Speed Model over time is clearly demonstrated in the simulated plots. Furthermore, the stark difference between the 1,000 and 10,000 realization graphs resolves a major implementation question from earlier discussions. It proves definitively that any observed gaps between the theoretical formulas and the empirical tests were artifacts of an insufficient sample size rather than a flaw in the theoretical framework or the duality equivalence.

% ------------------------------------------------
% SCRIBE QUESTION 6
% ------------------------------------------------

\section*{Scribe Question 6: Open Issues and Responsibility Attribution}

\textbf{Computational Bottlenecks and Boundary Constraints}

Although scaling the simulation to 10,000 realizations was able to smooth the empirical. - curves, this success throws into sharp relief an unresolved problem with respect to computational complexity. Operating tens of thousands of independent network instantiations across. spatial bins and many time steps cause gigantically large processing overhead. To reach the gold standard of 1,000,000 realizations that is commonly desired in rigorous stochastic geometry. sub research, the existing iterative cycles have a harsh computational bottleneck. Addition- ally, one has a hypothetical unresolved problem on the spatial limit of the simulation. The mathematical model assumes that there is an infinite two-dimensional plane but the code imposes a mathematical model that there is an infinite two-dimensional plane but the code imposes a finite 10-kilometer radius. Whether this hard boundary is an artificially created boundary is still yet to be verified and it artificially distorts the incoming density of drones with long simulation periods because of edge 

\textbf{Team Responsibilities and Next Steps}

In order to address these pending issues, one should delegate certain duties in the following stage. of the project. The programming head should be concerned with optimization of Monte Carlo execution, experimenting with more complex vectorization, or use of parallel processing. libraries in order to get extreme realization counts computationally feasible without crashing. the system. The mathematical analyst has the responsibility of exploring the finite boundary effects, proving mathematically whether the 10-kilometer limit is a perfect approximation of the. that is, infinite-plane assumption, and that the numerical integration functions do so as to preserve precision at edge limits. Lastly, the scribe has the duty of recording these in detail. certain computational tradeoffs, and there will be a final report that represents the work of the team. 

\end{document}